\subsubsection{Funktionsumfang}
Folgender Funktionsumfang wurde fuer das System definiert:
\begin{itemize}
\item Landing Page mit allgemeinen Infos zum Betrieb und Weiterleitungsmoeglichkeit zur Stellenuebersicht 
\item Eine Uebersicht aller offener Stellen mit individueller Weiterleitung zur angeklickten Stellenanzeige
\item Die Stellenanzeige mit allgemeinen Infos zur Stelle, Aufgabenbeschreibung, Vorraussetzungen und Weiterleitungsbutton zum Bewerbungsformular
\item Das Bewerbungsformular, wo die allgemeinen Kontaktdaten zu sehen sind, Bewerbungsunterlagen (Anschreiben, Lebenslauf, Zeugnisse, Zertifikate) hochgeladen werden muessen und die zu bewerbende Stelle ausgewaehlt werden muss
\item Das Bewerberdashboard, welches eine Statuseinsicht aller Bewerbungen darstellt
\item Die Admin-Panel-Uebersicht fuer den Recruiter. Hier wird es eine Uebersicht aller Bewerbungen geben. Es wird die Moeglichkeit bestehen in die Detailanzeige einer bestimmten Bewerbung zu wechseln.
\item Die Admin-Panel-Bewerbungsanzeige. Hier wird die in der Uebersicht ausgewaehlte Bewerbung detailiert angezeigt und es besteht hier die Moeglichkeit die hochgeladenen Unterlagen sich runterzuladen.
\item Die Admin-Panel-Stellenerstellung bzw. -editierung. Hier soll die Moeglichkeit fuer den Recruiter geschaffen werden, neue Stellen oder aktive Stellen zu bearbeien bzw. deaktivieren.
\item Ein Login-Fenster muss implementiert werden, sodass sowohl Bewerber und Recruiter auf die entsprechenden Dashboards zugreifen koennen.
\item Fuer neue Bewerber wird ebenfalls eine Register-Page eingerichtet, wo allgemeine Kontaktdaten sowie Mail und Passwort fuer kuenftige Anmeldungen vergeben werden.
\item Im Falle, dass ein Bewerber sein Passwort vergisst, wird ein Passwort-Reset Prozess eingebaut.
\end{itemize}
Dies sind die zentralen Funktionaliaeten, die als Basis fuer Architektur, UI und Datenbankmodellierung dienen.
